% Options for packages loaded elsewhere
\PassOptionsToPackage{unicode}{hyperref}
\PassOptionsToPackage{hyphens}{url}
\documentclass[
]{ltjsarticle}
\usepackage{xcolor}
\usepackage{amsmath,amssymb}
\setcounter{secnumdepth}{-\maxdimen} % remove section numbering
\usepackage{iftex}
\ifPDFTeX
  \usepackage[T1]{fontenc}
  \usepackage[utf8]{inputenc}
  \usepackage{textcomp} % provide euro and other symbols
\else % if luatex or xetex
  \usepackage{unicode-math} % this also loads fontspec
  \defaultfontfeatures{Scale=MatchLowercase}
  \defaultfontfeatures[\rmfamily]{Ligatures=TeX,Scale=1}
\fi
\usepackage{lmodern}
\ifPDFTeX\else
  % xetex/luatex font selection
\fi
% Use upquote if available, for straight quotes in verbatim environments
\IfFileExists{upquote.sty}{\usepackage{upquote}}{}
\IfFileExists{microtype.sty}{% use microtype if available
  \usepackage[]{microtype}
  \UseMicrotypeSet[protrusion]{basicmath} % disable protrusion for tt fonts
}{}
\makeatletter
\@ifundefined{KOMAClassName}{% if non-KOMA class
  \IfFileExists{parskip.sty}{%
    \usepackage{parskip}
  }{% else
    \setlength{\parindent}{0pt}
    \setlength{\parskip}{6pt plus 2pt minus 1pt}}
}{% if KOMA class
  \KOMAoptions{parskip=half}}
\makeatother
\setlength{\emergencystretch}{3em} % prevent overfull lines
\providecommand{\tightlist}{%
  \setlength{\itemsep}{0pt}\setlength{\parskip}{0pt}}
\usepackage{bookmark}
\IfFileExists{xurl.sty}{\usepackage{xurl}}{} % add URL line breaks if available
\urlstyle{same}
\hypersetup{
  hidelinks,
  pdfcreator={LaTeX via pandoc}}

\author{}
\date{}

\begin{document}

\section{あるモデル粒子をめぐる思考実験（続）------総和ゼロ条件から見る第五の振動数と合成曲率}\label{ux3042ux308bux30e2ux30c7ux30ebux7c92ux5b50ux3092ux3081ux3050ux308bux601dux8003ux5b9fux9a13ux7d9aux7dcfux548cux30bcux30edux6761ux4ef6ux304bux3089ux898bux308bux7b2cux4e94ux306eux632fux52d5ux6570ux3068ux5408ux6210ux66f2ux7387}

\textbf{著者}：木原 範昭（Noriaki Kihara） \textbf{所属}：WF System Co.,
Ltd.~/ 大阪大学基礎工学部（卒業）
\textbf{ORCID}：\href{https://orcid.org/0009-0004-6753-4020}{0009-0004-6753-4020}
\textbf{版}：v1（ステルス・ドラフト） \textbf{日付}：2026 年 6 月
\textbf{ライセンス}：CC BY 4.0 \textbf{Concept
DOI}：\href{https://doi.org/10.5281/zenodo.20535894}{10.5281/zenodo.20535894}
\textbf{Version DOI
(v1.0)}：\href{https://doi.org/10.5281/zenodo.20535895}{10.5281/zenodo.20535895}

\begin{center}\rule{0.5\linewidth}{0.5pt}\end{center}

\subsection{本稿の性格}\label{ux672cux7a3fux306eux6027ux683c}

\textbf{本稿は観察論文であり、思考実験の形式を採る。証明論文・主張論文ではない。}

本稿は前稿
{[}14{]}（あるモデル粒子をめぐる思考実験------位相面積の床から次元の許容条件へ）の続編である。前稿が、点になれない床（位相空間の面積ゼロ状態の不在）から、内在する振動自由度の本数が四に許容される条件までを観察したのに対し、本稿は、その四本にエネルギーが伴うとき、総和ゼロ型の保存条件を課すと第五の方向が形式上現れること、およびその第五の方向を半径とする中心投影と、続く軸での切断による合成曲率半径までを、同じ一つの思考実験の延長として観察する。

本稿は標準量子論・一般相対論・特殊相対論を修正しない。新しい物理理論を提案しない。新しい数学的定理を証明しない。観測可能な予言を一切変更しない。本稿が行うのは、既知の構造------四元運動量がエネルギーと運動量を区別しない一本のベクトルであること（特殊・一般相対論
{[}10{]}）、Robertson の不確定性の床
{[}2{]}、中心投影による次元削減の代数構造
{[}19{]}------を、振動数という一つの基本量の視点で並べ直して観察することに尽きる。

本稿は以下を主張しない：

\begin{itemize}
\tightlist
\item
  標準量子論・一般相対論・特殊相対論の数学的予言の変更
\item
  時空次元・計量符号・ミンコフスキー構造の物理的導出
\item
  質量・エネルギー・時間・運動量が物理的実体ではない、という存在論的主張（後述のとおり、本稿はこれらを記述上どう派生させるかを扱うのであって、実体性を否定しない）
\item
  第五の振動数を静止質量と同定すること（構造の一致を観察するに留める）
\item
  新しい物理定数・宇宙定数の値の予言
\item
  新しい数学的定理の証明
\item
  宇宙の起源・時間の本性に関する形而上学的主張
\end{itemize}

評価・解釈は読者に委ねる。

\begin{center}\rule{0.5\linewidth}{0.5pt}\end{center}

\subsection{基本計量についての宣言------本稿が何を基本に据えるか}\label{ux57faux672cux8a08ux91cfux306bux3064ux3044ux3066ux306eux5ba3ux8a00ux672cux7a3fux304cux4f55ux3092ux57faux672cux306bux636eux3048ux308bux304b}

本稿の全体を貫く一つの取り方を、最初に明示する。これは以降のすべての節の前提である。

\textbf{本稿は、記述の基本計量（基本自由度）を振動数 \(\nu\) に取る。}
時間・運動量・エネルギーといった量は、\(\nu\)
と独立な基本自由度としてではなく、\(\nu\)
から派生する従属量として扱う。これは記述上の取り方であって、時間や運動量やエネルギーが物理的実体ではないと主張するものではない。これらが物理的実体として機能することは標準物理のとおりであり、本稿はそれと矛盾しない。本稿の限定は、何を基本に据えて記述を組むかという選択にとどまる。

この取り方を採る理由は、時間という語が物理において多義的だからである。状態の発展を測るパラメータとしての時間と、計量における座標軸としての時間は、平坦時空・時間並進対称性のもとでは一致するが、一般にはその同一性が自明でない（一般相対論における「時間の問題」）。本稿が時間を基本語彙に置くと、この多義性が無意識に紛れ込み、証明していない同一視を証明済みのように扱う危険がある。これを避けるため、本稿は基本語彙を
\(\nu\) に統一し、時間に見える量は \(\nu\)
からの派生物としてのみ導入する。

この宣言、および本稿後半（§3
以降）の振動自由度から時空構造への接続は、前稿 {[}14{]}
と同じく\textbf{著者の対応仮説}であって、標準理論からの導出ではない。確立した数学（保存則・Robertson
の床 {[}2{]}・中心投影の合成演算
{[}19{]}・多元体の次元系列）を第一層、それらを振動自由度へ接続して読む部分（中心投影による定式化）を第二層（対応仮説）とする。以降はこの二層の区別のもとで読まれたい。

\begin{center}\rule{0.5\linewidth}{0.5pt}\end{center}

\subsection{要旨}\label{ux8981ux65e8}

前稿 {[}14{]}
は、内在する振動自由度の本数が四に許容されることを観察した。本稿はその四本にエネルギーが伴うときに何が起きるかを、思考実験の延長として追う。

第一に、四本の振動数が対等な段階では、エネルギーと運動量は区別されない。一般相対論・特殊相対論において四元運動量
\(p^\mu=(E,\boldsymbol{p})\)
は一本のベクトルの四成分であり、エネルギーと運動量の分離は観測者が時間軸を選んだ結果である。四本対等の段階では「エネルギー相当の量」は四元ベクトル全体（およびその不変量）として在るが、特定の一成分としては未分化である。

第二に、背景時空も絶対スケールも持たない系で、このエネルギーが正の側に在るとき、総和ゼロ型の保存条件を課すと、これを相殺する負の方向が形式上現れる（標準理論からの導出ではない）。四本の中に相殺相手はないため、第五の方向
\(\nu_5\) が負を担う。すると
\(\nu_1^2+\nu_2^2+\nu_3^2+\nu_4^2-\nu_5^2=0\) が得られる。これは署名
\((+,+,+,+,-)\) の光円錐ヌル条件であり、また四次元から見た質量殻
\(E^2=\boldsymbol{p}^2+m^2\) と同じ二次形式である（本稿は \(\nu_5\)
を質量と同定しない。標準の四次元質量殻に寄せて読むと、孤立軸 \(\nu_5\)
はエネルギー様の位置に、質量項の位置には四本のうち一本が現れるという、構造の一致を観察するに留める）。\(\nu_5\)
は独立自由度ではなく四本から決まる従属量（合成半径）であり、四本と同じく床
\(\tfrac12\) の不確定性を負う。

第三に、第五の方向 \(\nu_5\)
を半径とする中心投影を定式化する（第一段階）。各軸の値を振動数の逆数
\(1/\nu_i\) に取り、第五成分を \(R=1/\nu_5\) として、半径 \(R\)
の球面への中心投影を施す。続いて、別の軸での球面上の切断を定式化し（第二段階）、切断後の合成曲率半径が残存座標の二乗和
\(r_{\rm final}^2=R^2-\sum_{k\in S}x_k'^2\)
の閉形式で与えられることを示す。これらは中心投影の合成演算 {[}19{]}
に厳密に従う幾何・代数の計算であり、解釈・観測的主張は伴わない。

本稿は新しい主張を行わず、これらの構造を一つの思考実験の中で並置観察するに留まる。符号構造の起源、\(\nu_5\)
と静止質量の対応、投影・切断の物理的解釈は、今後の課題として明示的に外す。

\begin{center}\rule{0.5\linewidth}{0.5pt}\end{center}

\subsection{§1
出発点の確認------前稿の到達点}\label{ux51faux767aux70b9ux306eux78baux8a8dux524dux7a3fux306eux5230ux9054ux70b9}

前稿 {[}14{]}
の到達点を引き継ぐ。内在する振動自由度は、点になれない床（\(\sqrt{\det\Sigma}\ge\tfrac12\)
{[}4{]}、面積ゼロ状態の不在）と、軸の合成が閉じる離散条件（可除・非可換・結合を満たす最小次元、四元数
\(\mathbb{H}\)）のもとで、四本に許容される。四本の振動数
\(\nu_1,\nu_2,\nu_3,\nu_4\)
それぞれに、対応する従属量（運動量様の成分）がぶら下がる。

この結論を一段落で再掲する。Robertson の不確定性 {[}2{]}
は、位相空間に絞りきれない最小面積------床
\(\sqrt{\det\Sigma}\ge\tfrac12\)、de Gosson の量子ブロブ
{[}4{]}------を課す。この床のため状態は一点には潰れられず、内在的な振動自由度を持たざるをえない。さらに、その振動軸どうしの合成が閉じる（可除・非可換・結合を満たす）最小次元は、多元体の次元系列
\(1,2,4,8\) のうち四元数
\(\mathbb{H}\)（次元四）である。床による振動の不可避性と、\(\mathbb{H}\)
による軸合成の閉包------この二つから、内在する振動自由度の本数が四に許容される。本稿はこの対等な四本を出発点とする。

前稿はここで、四本の従属量を並べた「形式的な四成分量」を得たが、それが物理的な四元運動量になるにはミンコフスキー計量の別途導入が要るとして、符号構造を外していた。本稿はその外した部分の一歩手前------エネルギーが四成分のどこに、どう現れるか------から再開する。

なお前稿の §4.2 で観察したとおり、二軸の比 \(k=\nu_2/\nu_1\)
は等方（\(k=1\)）と非等方（\(k\neq1\)）を分ける。\(k=1\)
では空間方向と呼べる特権的な軸がまだ立っておらず、四本は対等である。本稿の出発点はこの対等な四本の段階である。

\begin{center}\rule{0.5\linewidth}{0.5pt}\end{center}

\subsection{§2
エネルギーと未分化の四元運動量}\label{ux30a8ux30cdux30ebux30aeux30fcux3068ux672aux5206ux5316ux306eux56dbux5143ux904bux52d5ux91cf}

\subsubsection{2.1
局所慣性系ではエネルギーと運動量は四元運動量の成分として統一される}\label{ux5c40ux6240ux6163ux6027ux7cfbux3067ux306fux30a8ux30cdux30ebux30aeux30fcux3068ux904bux52d5ux91cfux306fux56dbux5143ux904bux52d5ux91cfux306eux6210ux5206ux3068ux3057ux3066ux7d71ux4e00ux3055ux308cux308b}

特殊相対論および一般相対論の局所慣性系において、四元運動量
\(p^\mu=(E/c,\boldsymbol{p})\) は一本の四元ベクトルであり、エネルギー
\(E\) と運動量 \(\boldsymbol{p}\)
は別々の量ではなく、同じベクトルの四成分である。両者の分離は観測者に相対的である。四元速度
\(u^\mu\) を持つ観測者にとってのエネルギーは射影 \(E=p^\mu u_\mu\)
で与えられ、運動量はそれに直交する三成分である。観測者を変えれば（ブースト）、同じ
\(p^\mu\)
のエネルギー成分と運動量成分の振り分けが変わる。これは標準的な事実であり、本稿が新たに述べることではない。

\subsubsection{2.2
四本対等の段階で言えること}\label{ux56dbux672cux5bfeux7b49ux306eux6bb5ux968eux3067ux8a00ux3048ux308bux3053ux3068}

本稿の四本の振動数が四元運動量に対応するなら、どれが時間成分（エネルギー）かが未分化な段階では、言えることは限られる。すなわち、四つの対等な成分（四元運動量という一本のベクトル）が在ること、そしてそのベクトルの観測者によらない不変な大きさ
\(p^\mu p_\mu\) が在ること、の二つである。

したがって「エネルギー相当の量は存在するか」への答えは、二つの意味で肯定される。四元ベクトル全体として在り、かつ不変量として在る。ただし「エネルギー」を時間成分という特定の一本として言うなら、それは観測者の時間軸選択を待つ。前者の意味では四本対等の段階で在り、後者の意味ではまだ分化していない。本稿が以降「エネルギーがある」と言うとき、それは前者（四元ベクトルおよびその不変量）を指す。強調すれば、四本対等の段階では「エネルギー」という特定の一成分はまだ分化しておらず、その分化は
§3 以降の段を待つ。

ここで、基本計量の宣言が効く。エネルギーは時間方向の振動数の従属量であり、運動量は空間方向の振動数の従属量である------\(E=h\nu\)
と \(\boldsymbol{p}=h\boldsymbol{k}\)
は、別々の関係式ではなく、「振動数に従属量がつく」という同一構造の四成分である。エネルギーと運動量の区別が観測者の時間軸選択に相対的であることは、振動数のどの一本を時間方向と読むかが未分化であることの、別の言い方になる。

\begin{center}\rule{0.5\linewidth}{0.5pt}\end{center}

\subsection{§3
無から在る粒子と保存則}\label{ux7121ux304bux3089ux5728ux308bux7c92ux5b50ux3068ux4fddux5b58ux5247}

\subsubsection{3.1
総和ゼロの条件}\label{ux7dcfux548cux30bcux30edux306eux6761ux4ef6}

何もない空間に、一つのモデル粒子が在るとする。背景もエネルギーの供給源も外には無い（本稿の系には背景時空も絶対スケールも無い）。にもかかわらず、§2
のエネルギーが在る。

ここで、本稿の思考実験として、全体の二次形式がゼロになるという保存型条件を課す。これは時間並進対称性から導かれる
Noether
保存則ではなく、背景もスケールも持たない本系に対して形式的に課す条件である（標準理論からの導出ではない）。四本の振動数に伴う成分（四元運動量の成分）が正の側に在る以上、それを相殺する負の寄与がなければ、二次形式の総和はゼロにならない。

四本の中には相殺相手がない（四本は対等で同じ側に在る）。したがって、この保存型条件を満たすには、四本の外の第五の方向が負の寄与を担えばよい。第五の方向の寄与が負であれば、四本の正と打ち消し合い、二次形式の総和がゼロになる。これは標準理論からの導出ではなく、保存型条件を課したときに形式上現れる第五成分である------本稿はこの読みを思考実験として採る。

\subsubsection{3.2
第五の方向の負符号（便宜的に虚方向と呼ぶ）}\label{ux7b2cux4e94ux306eux65b9ux5411ux306eux8ca0ux7b26ux53f7ux4fbfux5b9cux7684ux306bux865aux65b9ux5411ux3068ux547cux3076}

この負の寄与を担う第五の方向は、観測される実空間の追加軸として現れるものではない。負の寄与は、二次形式の負符号------署名
\((+,+,+,+,-)\) の負号------として担われる。これを「\(\nu_5\)
が虚である」と書くこともできるが、それは計量符号を便宜的に言い換えた表現であって、\(\nu_5\)
を複素数とする（座標を複素化する）ことや、物理的に観測不能な実在方向を導出することを意味しない。本稿は負号を符号として扱い、第五の方向を便宜的に虚方向・R
方向と呼ぶにとどめる（曲率半径方向に対応する。§4.4--§5
参照）。この読みは、前稿 {[}14{]} および同著者 {[}15{]}
における「観測されない方向」の読みと、符号の上で整合する。

\begin{center}\rule{0.5\linewidth}{0.5pt}\end{center}

\subsection{\texorpdfstring{§4 第五の振動数
\(\nu_5\)}{§4 第五の振動数 \textbackslash nu\_5}}\label{ux7b2cux4e94ux306eux632fux52d5ux6570-nu_5}

\subsubsection{4.1
四平方和とヌル条件}\label{ux56dbux5e73ux65b9ux548cux3068ux30ccux30ebux6761ux4ef6}

ここで \(\nu_i\)
は、通常の周波数そのものではなく、四元運動量の各成分を周波数次元へ規格化した形式変数である（時間成分は
\(E\to E/h\)、空間成分は \(p\to pc/h\) のように、運動量側には \(c\)
が入る。基本計量の宣言）。総和ゼロの条件を、この形式変数の二乗で書くと

\[\nu_1^2+\nu_2^2+\nu_3^2+\nu_4^2-\nu_5^2=0\]

すなわち

\[\nu_5^2=\nu_1^2+\nu_2^2+\nu_3^2+\nu_4^2,\qquad \nu_5=\sqrt{\nu_1^2+\nu_2^2+\nu_3^2+\nu_4^2}\]

である。\(\nu_5\)
は独立な自由度ではなく、四本から決まる従属量である。これは本稿で課した保存型条件（§3）の言い換えであり、また同著者の合成半径
\(R=\sqrt{\sum R_k^2}\)（複数軸の球面投影の合成）と同じ形である。

この式は三つの顔を持つ。第一に、本稿で課した総和ゼロ型の保存条件（§3）である。第二に、署名
\((+,+,+,+,-)\)
の五次元ミンコフスキー空間における光円錐のヌル条件である（四本が同符号、\(\nu_5\)
が符号の異なる孤立軸）。第三に、これを四次元の側から見ると、四本のうち一本を凍結した像が、質量殻
\(E^2-\boldsymbol{p}^2=m^2\)
と（全体符号を除いて）同じ二次形式になる。ここで項ごとの対応を正確にとると、孤立軸
\(\nu_5\) は標準の質量殻の孤立項 \(E^2\) の位置に、質量項 \(m^2\)
の位置には凍結した四本のうちの一本が対応する（孤立軸 \(\nu_5\)
を質量項に置くと、\((4,1)\) の符号数のもとでは殻ではなく四変数の球面
\(S^3\) になり、対応が崩れる）。これは §4.2 の Kaluza--Klein
型の読み------四本のうちの一本が低次元で質量として現れる------と一致する。物理の三つの記述が、一つのヌル条件に畳まれている。

\textbf{符号構造の起源は外す。}
以上の三つの顔は、語に依存しない算術上の「形の一致」の観察である。どの軸が質量・エネルギー・時間に当たるかという指示対象の割り当ては、背景を持たない本系では未分化であり、本稿は
commit しない（上の対応は、標準の四次元質量殻へ寄せた as-if
の位置対応にすぎない）。また、四本のうち一本が他の三本に対して負符号を獲得して通常の質量殻の符号
\((+,-,-,-)\)
へ移行する機構（ユークリッドからミンコフスキーへの符号構造の起源）は、本稿では導出せず、今後の課題として外す。

\textbf{主張の限定。} 本節が \(\nu^2\)
をエネルギー・運動量と結ぶのは、二次形式の構造的な一致------同じ形が現れること------の観察であって、\(\nu^2\)
がエネルギー・運動量の物理的実体であると主張するものではない。「なぜ
\(\nu^2\)
がエネルギーなのか」という問いには本稿は立ち入らない。観察するのは、保存条件・光円錐ヌル条件・質量殻が一つの二次形式に畳まれるという、形の一致だけである（基本計量の宣言に従い、エネルギー・運動量は
\(\nu\) からの派生従属量として扱う）。

\subsubsection{4.2
質量との同定はしない}\label{ux8ceaux91cfux3068ux306eux540cux5b9aux306fux3057ux306aux3044}

第二・第三の顔について、正確に述べる。\(\nu_5^2=\sum\nu_k^2\)
が、五次元では質量ゼロのヌル条件であり、四次元から見ると質量殻と同じ形になることは、構造の一致である。高次元のヌル（質量ゼロ）が、余剰方向の運動量ぶんだけ低次元で質量として現れる、という対応（Kaluza--Klein
型 {[}16{]} の機構と同じ形）が読める。

しかし本稿は、\(\nu_5\) を静止質量と同定しない。§4.1
で整理したとおり、署名 \((+,+,+,+,-)\)
の五次元ヌル条件を標準的な四次元質量殻の形に寄せて読むと、孤立符号を持つ
\(\nu_5\)
は孤立項（エネルギー様）の位置に現れ、質量項の位置には四本のうち凍結・切断された一本が現れる。したがって本稿が観察するのは、\(\nu_5\)
そのものが質量であるという対応ではなく、高次元のヌル条件を低次元側から読むと質量殻型の二次形式が現れる、という構造の一致である。\(\nu_5\)
が質量であると主張するには、\(\nu_5\)
が質量の他の性質（慣性・重力源・ローレンツ不変な静止量）も満たすことの確認が要るが、本稿はその確認を行わず、構造の一致の観察に留める。

\subsubsection{\texorpdfstring{4.3 \(\nu_5\)
にも床がある}{4.3 \textbackslash nu\_5 にも床がある}}\label{nu_5-ux306bux3082ux5e8aux304cux3042ux308b}

\(\nu_5\) は四本から決まる従属量である。四本それぞれが床 \(\tfrac12\)
を負って確定した整数に固定できない以上、その二乗和の平方根である
\(\nu_5\)
もまた絞りきれない半端を負い、確定した整数には固定できない。すなわち
\(\nu_5\)
の不確定性は、四本のゆらぎが二乗和の平方根を通じて伝播してくる分として尽き、\(\nu_5\)
に独立な床を別途加えるものではない（独立自由度の追加を意味しない）。前稿
{[}14{]} の「整数 \(n\) ＋絞りきれない半端
\(\tfrac12\)」という構造が、この伝播を通じて第五の方向にも及ぶ。

\subsubsection{\texorpdfstring{4.4 R 軸（\(\nu_5\)
方向）中心投影の第一段階}{4.4 R 軸（\textbackslash nu\_5 方向）中心投影の第一段階}}\label{r-ux8ef8nu_5-ux65b9ux5411ux4e2dux5fc3ux6295ux5f71ux306eux7b2cux4e00ux6bb5ux968e}

第五の方向 \(\nu_5\) を半径とする中心投影の第一段階を定式化する。本節は
R 軸での 1 回の中心投影のみを扱い、他軸（\(t=1/\nu_4\)
等）での切断は行わない。

\textbf{定義についての断り。}
本節では、各軸の座標値を対応する振動数の逆数 \(x_i:=1/\nu_i\)
と\textbf{定義する}。これは中心投影を幾何学的に書き下すための座標の取り方（記法上の定義）に過ぎず、\(x_i=1/\nu_i\)
が物理的な実像（長さ・波長などの実在量）であると主張するものではない。本節以降の操作は、すべてこの定義のもとでの幾何・代数の計算である。

\textbf{座標。} 各軸の値を対応する振動数の逆数に取り、第五成分を
\(R=1/\nu_5\) とする 5 次元の点を

\[x = \left(\frac{1}{\nu_1},\ \frac{1}{\nu_2},\ \frac{1}{\nu_3},\ \frac{1}{\nu_4},\ R\right),\qquad R=\frac{1}{\nu_5}\]

とする。ここで \(\nu_5\) は §4.1 の保存型条件
\(\nu_1^2+\nu_2^2+\nu_3^2+\nu_4^2=\nu_5^2\) を満たす。

\textbf{ノルム。}

\[|x| = \sqrt{\frac{1}{\nu_1^2}+\frac{1}{\nu_2^2}+\frac{1}{\nu_3^2}+\frac{1}{\nu_4^2}+R^2} = \sqrt{\sum_{i=1}^{4}\frac{1}{\nu_i^2}+\frac{1}{\nu_5^2}}\]

\textbf{中心投影。} 原点を中心とし半径 \(R\) の球面への中心投影

\[x' = \frac{R}{|x|}\,x\]

を施す。スケール因子を \(\lambda := R/|x|\) と置くと、各成分は

\[x_i' = \lambda\,x_i = \frac{R}{|x|}\cdot\frac{1}{\nu_i}\quad(i=1,2,3,4),\qquad R' = \lambda\,R = \frac{R^2}{|x|}\]

\textbf{投影後の振動数。} \(\nu_i'=1/x_i'\)、\(\nu_5'=1/R'\) により

\[\nu_i' = \frac{|x|}{R}\,\nu_i\quad(i=1,2,3,4),\qquad \nu_5' = \frac{1}{R'} = \frac{|x|}{R^2} = \frac{|x|}{R}\,\nu_5\]

すなわち \(i=1,\dots,5\) のすべてについて

\[\nu_i' = \frac{|x|}{R}\,\nu_i\]

である。スケール因子 \(|x|/R\) を \(\nu\) で書けば
\(|x|/R = |x|\,\nu_5 = \sqrt{1+\nu_5^2\sum_{j=1}^{4}\nu_j^{-2}}\)
であり、したがって

\[\nu_i' = \nu_i\sqrt{1+\nu_5^2\sum_{j=1}^{4}\frac{1}{\nu_j^2}}\qquad(i=1,\dots,5)\]

として、投影後の点
\(x'=(1/\nu_1',\,1/\nu_2',\,1/\nu_3',\,1/\nu_4',\,R')\)
の各成分が定式化される。

\textbf{\(R'\) の検算（R 軸上の点）。} 投影の中心軸上の点
\(x=(0,0,0,0,R)\) では \(|x|=R\) であり、

\[x' = \frac{R}{|x|}\,x = \frac{R}{R}\,x = x,\qquad R' = \frac{R^2}{|x|} = \frac{R^2}{R} = R\]

となる。すなわち R 軸上では \(R'=R\)
で、半径成分は投影によって変化しない。

\begin{center}\rule{0.5\linewidth}{0.5pt}\end{center}

\subsection{§5
任意軸での投影------第二段階の切断}\label{ux4efbux610fux8ef8ux3067ux306eux6295ux5f71ux7b2cux4e8cux6bb5ux968eux306eux5207ux65ad}

本節は、§4.4 の第一段階（R
軸中心投影）に続く第二段階を定式化する。中心投影は 1
回限りであり、続く軸での操作は球面上の切断 \(\sigma\)
である（{[}19{]}）。具体例として軸 \(x_4=1/\nu_4\)
での切断を計算し、任意軸への一般形を与える。本節は幾何・代数の計算に限る。

\textbf{出発点。} §4.4 の第一段階の後、点は半径 \(R=1/\nu_5\) の球面
\(S^4(R)\) 上にある：

\[x'=\frac{R}{|x|}\,x=(x_1',x_2',x_3',x_4',x_5'),\qquad |x'|=R\]

ここで
\(x_i'=\dfrac{R}{|x|}\cdot\dfrac{1}{\nu_i}\)（\(i=1,2,3,4\)）、\(x_5'=R'=\dfrac{R^2}{|x|}\)
である。球面上であることから

\[x_1'^2+x_2'^2+x_3'^2+x_4'^2+x_5'^2=R^2 \tag{5.1}\]

\textbf{軸 \(x_4\) での切断 \(\sigma_{x_4}\)。} 球面 \(S^4(R)\) 上の点の
\(x_4\) 成分を、その球面上の値 \(c=x_4'\) に固定する。超平面
\(\{x_4=c\}\) と \(S^4(R)\) の交わりは、\(x_4\) 軸を除いた半径 \(r''\)
の球面 \(S^3(r'')\) に同型であり、

\[r''=\sqrt{R^2-c^2}=\sqrt{R^2-x_4'^{\,2}}\]

切断によって除かれる \(x_4\) 軸以外の残存座標 \(x_1',x_2',x_3',x_5'\)
は変化しない。

\textbf{閉形式。} (5.1) より
\(R^2-x_4'^{\,2}=x_1'^2+x_2'^2+x_3'^2+x_5'^2\) であるから、

\[r''^{\,2}=R^2-x_4'^{\,2}=x_1'^2+x_2'^2+x_3'^2+x_5'^2 \tag{5.2}\]

すなわち、切断後の半径は残存座標の二乗和の平方根である。

\textbf{振動数による表示。}
\(x_4'=\dfrac{R}{|x|\,\nu_4}=\dfrac{1}{\nu_4'}\)（\(\nu_4'=\dfrac{|x|}{R}\nu_4\)、§4.4）により

\[r''^{\,2}=R^2-\frac{R^2}{|x|^2\nu_4^2}=R^2\left(1-\frac{1}{|x|^2\nu_4^2}\right)\]

残存軸の振動数 \(\nu_1'',\nu_2'',\nu_3''\) は第一段階の値
\(\nu_1',\nu_2',\nu_3'\) に等しく（\(x_5'=R'\) も不変）、\(\nu_4\)
軸は切断で除かれる。

\textbf{任意軸・多軸への一般形。} 任意の軸
\(x_k\)（\(k\in\{1,2,3,4\}\)）での切断も同形で、

\[r''^{\,2}=R^2-x_k'^{\,2}=\sum_{i\neq k}x_i'^{\,2}\]

複数軸の集合 \(S\) で順次切断すれば、順序によらず

\[r_{\rm final}^{\,2}=R^2-\sum_{k\in S}x_k'^{\,2} \tag{5.3}\]

となる（{[}19{]}
定理2）。これが、任意軸での投影（第一段階＝中心投影、第二段階以降＝切断）の計算方法である。なお
\(r_{\rm final}\)
は、幾何学的には切断後の球面半径であり、本稿ではこれを（対応仮説の層で）合成曲率半径と呼ぶにとどめる。物理的・時空的な曲率を主張するものではない。

\begin{center}\rule{0.5\linewidth}{0.5pt}\end{center}

\subsection{参考文献}\label{ux53c2ux8003ux6587ux732e}

{[}2{]} H. P. Robertson (1929). \emph{The uncertainty principle}. Phys.
Rev.~\textbf{34}, 163. {[}4{]} M. de Gosson (2013). \emph{Quantum
blobs}. Found. Phys. \textbf{43}, 440. {[}10{]} C. W. Misner, K. S.
Thorne, J. A. Wheeler (1973). \emph{Gravitation}.
Freeman.（四元運動量、エネルギーと運動量の観測者相対性） {[}14{]} N.
Kihara (2026).
\emph{あるモデル粒子をめぐる思考実験------位相面積の床から次元の許容条件へ}.
観察論文（同著者・前稿）. {[}15{]} N. Kihara (2026). 複素双曲構造
\(\mathrm{CH}^2\)・\(\mathrm{SU}(2,1)\)
および合成曲率に関する同著者の論文. {[}16{]} T. Kaluza (1921); O. Klein
(1926).（高次元のヌルが低次元で質量として現れる機構の形） {[}19{]} N.
Kihara (2026). \emph{中心投影の合成演算と合成曲率半径の閉形式}.
観察論文（同著者）. Concept DOI: 10.5281/zenodo.20060728.

（注：文献番号は前稿 {[}14{]} の番号体系と整合させる。{[}14{]}{[}15{]}
の正式な DOI・書誌は公開時に確定する。）

\begin{center}\rule{0.5\linewidth}{0.5pt}\end{center}

著者：木原 範昭（Noriaki Kihara）／ WF System Co., Ltd.~／ ORCID
\href{https://orcid.org/0009-0004-6753-4020}{0009-0004-6753-4020}／ CC
BY 4.0

\begin{center}\rule{0.5\linewidth}{0.5pt}\end{center}

\subsection{改訂履歴}\label{ux6539ux8a02ux5c65ux6b74}

\begin{itemize}
\tightlist
\item
  \textbf{v1（2026-06）}：初版（ステルス・ドラフト）。前稿 {[}14{]}
  の続編。冒頭に基本計量の宣言（基本自由度を振動数 \(\nu\)
  に取り、時間・運動量・エネルギーを \(\nu\)
  からの派生従属量として扱う。ただし物理的実体性は否定せず、標準物理と矛盾しない）を置く。(§1)
  前稿の到達点（対等な四本の振動数）から再開。(§2)
  エネルギーと運動量は四元ベクトルの成分として未分化、分離は観測者相対的（一般相対論の標準
  {[}10{]}）。(§3)
  背景もスケールもない系で総和ゼロ型の保存条件を課すと第五の負の方向が形式上現れること、それは符号の上で虚方向（計量符号の負号）。(§4)
  \(\nu_1^2+\nu_2^2+\nu_3^2+\nu_4^2-\nu_5^2=0\)
  が保存条件・光円錐ヌル条件・質量殻型二次形式の三つの顔を持つこと。孤立軸
  \(\nu_5\) はエネルギー様の位置、質量は四本の一本。\(\nu_5\)
  は四本から決まる従属量（合成半径）、質量との同定はせず、\(\nu_5\)
  にも床 \(\tfrac12\)。(§4.4) 第五の方向 \(\nu_5\)
  を半径とする中心投影の第一段階を、座標
  \(x=(1/\nu_1,\dots,1/\nu_4,R)\)、\(R=1/\nu_5\) で定式化し、R 軸上で
  \(R'=R\) を検算。(§5) 続く軸 \(x_4=1/\nu_4\)
  での球面上の切断（第二段階）を定式化し、切断後の半径が残存座標の二乗和
  \(r_{\rm final}^2=R^2-\sum_{k\in S}x_k'^2\)
  の閉形式で与えられること（中心投影の合成演算 {[}19{]}
  と整合）を示す。二層（確立した数学／対応仮説）の区別を冒頭で宣言。本版で旧
  §5（自己参照・定常波の安定性）・旧 §6（合成曲率）・旧
  §7（時間の多義性）・旧 §8 を削除し、§4.4 と §5
  の幾何・代数の定式化に絞った。
\end{itemize}

\end{document}
